% ============================================================
%  FRONT MATTER
%  Title page · Copyright · Preface (pp. 3–6)
%  Statement of the Problem (pp. 7–14)
%  Manual Table of Contents
% ============================================================

% ---- Half-title page ----------------------------------------
\thispagestyle{empty}
\vspace*{6em}
\begin{center}
  {\large\bfseries\MakeUppercase{Spatial Constructions in Painting}}
\end{center}
\clearpage

% ---- Disclaimer of this translation page ---------------------------------------
\thispagestyle{empty}
\vspace*{6em}
\begin{center}
  {\large\bfseries\MakeUppercase{Disclaimer for this edition}} \par \bigskip
This document is an informal, AI-assisted English translation of \textrussian{Раушенбах Б.В. Пространственные построения в живописи. - М., 1980} [\textit{Prostranstvennye Kompozitsii v Zhivopisi (Moscow, 1980)} by Boris V. Rauschenbach], produced for personal reference. \par \smallskip
Translation was performed by Claude (Anthropic). Original figures, and plates are represented by placeholder captions only. Source pagination is preserved inline. This translation has not been reviewed for accuracy and should not be cited or treated as an authoritative rendering. All rights in the original work belong to their respective holders. \par  \bigskip

\textit{Project Page} \par
\href{https://github.com/star-lings/Rauschenbach-B-V-Spatial-constructions-in-painting-1980-English-Reproduction}{https://github.com/star-lings/Rauschenbach-B-V-Spatial-constructions-in-painting-1980-English-Reproduction} \par

\vfill

\textit{Version \today}

\end{center}

\begin{center}
    
\end{center}
\clearpage

% ---- Title page ---------------------------------------------
\thispagestyle{empty}
\vspace*{3em}
\begin{center}
  {\large Boris Viktorovich Rauschenbach}\\[2em]
  {\LARGE\bfseries SPATIAL CONSTRUCTIONS\\[0.4em] IN PAINTING}\\[3em]
  {\normalsize\itshape
    Based on the psychology of visual perception and with the aid\\
    of mathematical apparatus, the book analyses the theory of\\
    spatial constructions in the visual arts. Theoretical questions\\
    are examined through the material of art from Ancient Egypt,\\
    the Middle Ages (Byzantium, Ancient Rus\kern0pt', India, Iran)\\
    and the work of Cézanne. Mathematical derivations are set\\
    apart in Appendices. The book is illustrated.}\\[3em]
  {\normalsize Responsible editor: B.\,N.~Prokofiev}\\[4em]
  \vfill
  {\normalsize Nauka Publishers · Moscow · 1980}
\end{center}
\clearpage

% ---- Copyright / imprint page -------------------------------
\thispagestyle{empty}
\vspace*{\fill}
\noindent
{\small
Editor: F.\,I.~Grinberg. \quad
Artist: N.\,V.~Illarionova. \quad
Art editor: T.\,P.~Polenova.\\
Technical editor: R.\,M.~Denisova. \quad
Proof-readers: Yu.\,L.~Kosorygín, V.\,G.~Petrova.

\medskip\noindent
Registration No.~15142.\\
Submitted for composition 08.08.79. Signed for printing 07.12.79.\\
Format 70\,×\,90\,\textsuperscript{1}/\textsubscript{16}.
Print run 19\,400 copies. Order No.~2226. Price 2\,r.\ 10\,k.

\medskip\noindent
Nauka Publishers, GSP-7, Moscow, V-485, Profsoyuznaya ul., 90.\\
2nd Printing House of Nauka Publishers, 121099, Moscow, G-99,
Shubinsky per., 10.

\bigskip\noindent
\textcopyright\ Nauka Publishers, 1980
}
\vspace*{\fill}
\clearpage

% ============================================================
%  MANUAL TABLE OF CONTENTS
%  (Reflecting source pagination)
% ============================================================
\chapter*{Contents}
\addcontentsline{toc}{chapter}{Contents}
\thispagestyle{plain}

\begin{center}
  \begin{tabular}{@{}p{0.80\linewidth}r@{}}
    Preface  &  3 \\[0.4em]
    Statement of the Problem: Technical Drawing and Pictorial Rendering  &  7 \\[0.4em]
    \textit{Chapter\,I.}\ \ Painting and Relief of Ancient Egypt.
      Artistic Technical Drawing  &  15 \\[0.4em]
    \textit{Chapter\,II.}\ \ Preliminary Notes on Visual Perception  &  42 \\[0.4em]
    \textit{Chapter\,III.}\ \ Representing Visual Perception on a Plane  &  52 \\[0.4em]
    \textit{Chapter\,IV.}\ \ The Perspective Basis of Byzantine
      and Old Russian Painting  &  80 \\[0.4em]
    \textit{Chapter\,V.}\ \ Reverse Perspective  &  94 \\[0.4em]
    \textit{Chapter\,VI.}\ \ Effects of Reverse Perspective  &  118 \\[0.4em]
    \textit{Chapter\,VII.}\ \ Drafting Methods in Byzantine
      and Old Russian Painting  &  140 \\[0.4em]
    \textit{Chapter\,VIII.}\ \ Geometrically Contradictory Images  &  168 \\[0.4em]
    \textit{Chapter\,IX.}\ \ Indian and Iranian Miniatures.
      The System of Free Perceptual Perspective  &  190 \\[0.4em]
    \textit{Chapter\,X.}\ \ Cézanne's Space and the System of
      Rigid Perceptual Perspective  &  215 \\[0.4em]
    Conclusion  &  233 \\[0.4em]
    List of Abbreviations  &  242 \\[1em] 
      \end{tabular}
  \begin{tabular}{@{}p{0.80\linewidth}r@{}}{\textsc{Sketches of the Theory of Spatial Constructions in the Visual Arts}} \\[0.4em]
    \textit{Appendix\,1.}\ \ The Retinal Image and the System
      of Linear Perspective  &  244 \\[0.4em]
    \textit{Appendix\,2.}\ \ The Possibility of Perceptual Representation
      of Space on a Plane  &  247 \\[0.4em]
    \textit{Appendix\,3.}\ \ Quantitative Relationship Between Linear
      and Perceptual Perspectives  &  254 \\[0.4em]
    \textit{Appendix\,4.}\ \ Axonometry  &  258 \\[0.4em]
    \textit{Appendix\,5.}\ \ The Emergence of Reverse Perspective
      in Observing Near Regions of Space Due to
      the "Superconstancy" Phenomenon  &  260 \\[0.4em]
    \textit{Appendix\,6.}\ \ The Emergence of Reverse Perspective
      in Observing Objects in Foreshortening  &  263 \\[0.4em]
    \textit{Appendix\,7.}\ \ Reverse Perspective and Lobachevsky Geometry  &  271 \\[0.4em]
    \textit{Appendix\,8.}\ \ The Rigid System of Perceptual Perspective  &  275 \\[0.4em]
    \textit{Appendix\,9.}\ \ Zero Lines and Their Curvature  &  282 \\[0.4em]
    \textit{Appendix\,10.}\ \ Intuitive Representation of Sections
      of Four-Dimensional Space  &  285 \\[0.4em]
    \textit{Appendix\,11.}\ \ The Mechanism of Form Constancy
      and Portrait Painting  &  286 \\[0.4em]
  \end{tabular}
\end{center}
\clearpage

% ============================================================
%  PREFACE  [source pp. 3–6]
% ============================================================
\chapter*{Preface}
\addcontentsline{toc}{chapter}{Preface}
\markboth{Preface}{Preface}

The present book is devoted to the problem of the geometry of spatial constructions in
the visual arts. The term "spatial constructions" requires some clarification. As is well
known, when conveying a three-dimensional world on the flat surface of a painting, the
artist faces the task of representing the volumetric-plastic qualities of the depicted objects
and the task of constructing the space of the painting. Both tasks are not reducible to
geometry alone; no less important are light-and-shadow, colour, and the
spiritual-cultural experience organically fused with them, which together in synthesis
produce what is called "space in painting." The isolation of geometric features is always
conditional and always narrows the problem of space in painting, yet it proves useful and
allows recourse to the mathematical apparatus of research.

The content of the book is the geometric aspect of spatial constructions, encompassing
both the methods for depicting individual objects and the geometric methods for
constructing space as a whole — problems often subsumed under the concept of
"perspective systems."

The task to be studied is posed as follows: given (1) the laws of perception and (2)
geometry, what should the spatial constructions employed by the artist be, in order to
depict on a flat surface the actually perceived space as undistorted as possible? This
formulation of the problem allows the application of rigorous mathematical methods of
investigation. Of course, the artist is by no means obliged to slavishly follow nature; his
tasks are much broader. However, an understanding of how space ought to be depicted
according to the rules of geometry will allow the art historian to perceive more clearly
the methods and devices of the artist, to give a more subtle analysis of the work of
individual masters, and to arrive at a deeper characterization of the artistic peculiarities
of entire epochs.

The use of achievements of modern psychology of visual perception, in conjunction with
methods of mathematical analysis, has revealed the existence of a more complete system
of scientific perspective than had previously been recognized (called below "perceptual
perspective"), of which the linear perspective system of the Renaissance is a special case,
valid only for the isolated depiction of distant regions of space. This more complete
system "rehabilitates" (if such rehabilitation is needed at all) such peculiarities of
painting as the reverse perspective of the Middle Ages or Cézanne's space~--- not only
in connection with the distinctiveness of the artistic image, but also from the standpoint
of the strict logic of geometry. Moreover, it was discovered that in many cases the artist
by no means sets himself the goal of perspectival depiction of visible space, but considers
it necessary to depict the geometric properties of objective space. This task, too, was
examined from the standpoint of mathematics. Such an examination led to the conclusion
of the astounding perfection of the pictorial means employed by artists.

It should be emphasized that the author has everywhere approached works of painting
as examples illustrating these or those geometric properties of images, and in no way
undertakes to judge their artistic peculiarities, merits, or shortcomings. If the author
does give assessments and speaks of progress or regression in the historical development
of art, this is only from the point of view of compliance with, or violation of, strict
geometric logic~--- which is, of course, completely insufficient for the analysis of the
artistic image.

\srcpage{4}

The author has found it expedient to correlate the most typical geometric regularities
discovered in the course of mathematical analysis with the corresponding pictorial
material. In identifying such typical systems of geometric constructions, it emerged that
(excluding the most recent art from consideration) history knows in all only four basic
methods of spatial constructions on the plane of the image. These are: \textit{drafting
methods} (depiction of objective space, characteristic of, for example, the art of Ancient
Egypt), \textit{the method of local axonometries and their transformations} (corresponding
to ancient and medieval art), \textit{central linear perspective} of the Renaissance, and
\textit{central curvilinear perspective}, which appeared at the turn of the nineteenth and
twentieth centuries. Each of these four basic types of spatial constructions could manifest
itself at different times, in different regions, and in various modifications, giving rise to
an enormous diversity of pictorial means. This diversity was further enhanced by the fact
that the above-mentioned basic types of spatial constructions were frequently combined
with one another in the artistic works of any given culture. The study of all this
diversity~--- of the processes of sequential replacement of one method of spatial
construction by another, the causes that led to such changes, the problems of mutual
influence of artistic cultures, and other analogous questions~--- is not the content of the
present book. That is the domain of art history.

The basic content of the present book is, as already stated, the study of the fundamental
possibilities for conveying spatial formations on the plane of the image, proceeding from
the objective regularities of spatial perception inherent in human beings. In order that
the dry formal results obtained on this path might acquire vividness and persuasiveness,
they are applied to the analysis of the peculiarities of spatial constructions in the painting
of those epochs and regions where these peculiarities most fully and clearly reflect the
results obtained through mathematical investigation. It goes without saying that an
analogous analysis could be applied to other painting not discussed in the present book;
however, it seems that the fundamental regularities discovered probably exhaust the
methods of conveying space on a plane in the visual arts, and they are sufficiently fully
reflected in the material here presented for the reader's attention. This explains the
seemingly strange choice of topics: the book discusses the painting of Ancient Egypt, the
Middle Ages (Byzantium, Rus\kern0pt', India, Iran), and the work of Cézanne. Strictly
speaking, one ought also to consider the painting born of the Renaissance, but this is not
done, since its geometric structure is generally known. The primary attention in the book
is devoted to medieval art. This is explained by the fact that its geometric structure is the
most complex, that it involves and productively employs various approaches and methods,
and that in a certain sense this is geometrically the most diverse art.

\srcpage{5}

The present book arose as a development of the ideas set forth in the author's previous
work devoted to Old Russian painting.\footnote{%
  B.\,V.~Rauschenbach, \textit{Spatial Constructions in Old Russian Painting}. Moscow,
  1975. The expansion of the scope of application of the theory of spatial constructions
  was stimulated by oral and written responses received by the author, and in particular
  by the article of V.\,N.~Prokofiev, ``On `Perceptual Perspective' and Perspectives in
  Painting,'' included in that book.}
The treatment of problems connected with Old Russian painting has been supplemented
and refined; the sections devoted to the theory of two systems of perspective have been
substantially revised, and a rigorous exposition of the theory of reverse perspective has
been given. Taking into account numerous requests from readers, some facts and
considerations previously contained in the Appendices have been transferred to the main
text. Mathematical material has been gathered in the section "Sketches of the Theory of
Spatial Constructions in the Visual Arts." The main text of the book is generally
accessible; explanatory examples sufficient for understanding the substance of the
problems discussed are provided, and reference to the mathematical material is not at all
necessary. However, the mathematical sections cannot be regarded as inessential notes to
the main text, for they constitute the rigorous foundation on which the author's
argumentation rests. The effort to make this section more accessible to the reader
compelled the author to employ a mathematically elementary exposition presupposing
knowledge of mathematics at the level of a modern engineer's training. With the same
aim, the Appendices composing this section are arranged so that they may be read
independently of the main text of the book, providing in their main body (Appendices
1--9) a coherent exposition of the subject. The indicated section is accordingly divided
not into chapters but into Appendices, also because they can be read independently of
one another, as mathematical justifications for the corresponding sections of the main
text.

It should be noted that the present book solves the "direct" problem: how space should
be depicted on a plane under the condition of the closest possible fidelity to nature. In
principle there also exists the "inverse" problem: how will the viewer perceive this or
that arbitrary spatial construction created by the artist on the canvas? What spatial
representations will arise in the viewer? The difficulty of solving this problem is
connected not only with the fact that "inverse" problems are almost always more complex
than "direct" ones, but also with the fact that the perception of a work of art is not a
function of visual perception psychology and geometry alone. To no lesser~--- and rather
to a greater~--- degree it is a function of a person's experience of perceiving works of
art. But this latter circumstance takes the problem out of the realm of natural-scientific
conceptions into the domain of art history proper.

\srcpage{6}

In conclusion, the author considers it his pleasant duty to express gratitude to
A.\,I.~Komech and V.\,N.~Prokofiev, who undertook the labour of reading the manuscript
of the book and who offered valuable observations that were taken into account in its
final editing, as well as to Yu.\,G.~Gurevich, who carried out a careful editing of the
mathematical section of the book.

% ============================================================
%  STATEMENT OF THE PROBLEM: TECHNICAL DRAWING AND PICTORIAL
%  RENDERING  [source pp. 7–14]
% ============================================================
\chapter*{Statement of the Problem:\newline Technical Drawing and Pictorial Rendering}
\addcontentsline{toc}{chapter}{Statement of the Problem: Technical Drawing and Pictorial Rendering}
\markboth{Statement of the Problem}{Statement of the Problem}

\srcpage{7}

Wishing to analyse the fundamental possibilities available to the artist when depicting
three-dimensional objects, let us introduce the assumption that he is required to convey
their geometric features with documentary precision. It is perfectly clear that such a task
is never set before an artist; but this abstract formulation of the question is useful in the
present case, because the aim of the following discussion is precisely the analysis of the
fundamental possibilities available to the artist~--- and not the altogether different
question of the best method for conveying some artistic image.

Let us imagine that an artist undertakes a distant sea voyage and, during his travels, is
given the opportunity to see a crystal of the rarest beauty of form. How is he to convey,
to those at home, the form of this crystal? The artist who works in the domain of
three-dimensional sculpture (the maker of sculpture in the round) and the artist who
works in the domain of two-dimensional art (painter, graphic artist, maker of sculptural
reliefs, and so on) will find themselves in entirely different conditions when solving this
task. The former can in the end make a geometric copy of the crystal with any required
degree of precision and in natural size. In a completely different position is the latter,
confronted with the task of conveying three-dimensional space on a flat surface having
only two dimensions.

Let us call the first artist here a sculptor (without emphasizing each time that we are
speaking only of sculpture in the round) and the second a painter. The difficulty
confronting the painter is connected with the fact that man deals with two different
geometries of objects: first, the \textit{true geometry} of the object, that which it
possesses objectively; and second, the \textit{apparent geometry}, that which is
generated by the objective geometric form of the object but transformed in the course of
its visual perception. Speaking of this latter circumstance, artists usually assert that the
visible form of an object changes depending on the angle of view. The objective form of
an object can be apprehended by examining it from all sides, and still more by touching
it. Apparent form always presupposes purely visual perception, and this form will not
change only when the viewpoint on the object and the object itself are both stationary.

\srcpage{8}

It is important to emphasize that the geometry of objective space in the vast majority
of cases has no directly and immediately perceptible form. We see the surface of a
rectangular table in the form of a figure resembling a trapezoid or a parallelogram; that
it is in fact rectangular we know, but do not see.

In addition to objective space, man deals also with \textit{subjective} space, which is
better called \textit{perceptual} (from the word "perception," a synonym for the word
"sensation"). This directly visually perceivable space~--- the space of perception~--- is
characterized by a geometry sharply different from the geometry of objective space.
Everyone knows well that rails which are actually parallel appear to converge on the
horizon to a point.

When we speak of depicting space on a plane, we must first agree on which of these
two spaces is being discussed. If a flat image is required to convey the geometry of
objective space, let us agree to call such an image a \textit{technical drawing}, by no
means assuming it to be identical with modern engineering drawings (though these, too,
are intended for depicting the geometry of objective space) and fully recognizing all the
unusualness of this term in art-historical analysis. If a flat image is required to convey
the geometry of perceptual space, let us agree to call such an image a \textit{picture},
thereby emphasizing two things: that what is to be depicted is purely visual perception,
and that for the subsequent analysis only the geometric properties of such an image will
be essential, without regard to colour and other factors.

It is evident that a technical drawing can be obtained only as a result of the analytical
work of the mind (the comparison of diverse information about the depicted space),
while a picture is obtained by transferring one's visual perception onto the picture plane.

One must always clearly understand which of the two spaces the artist is depicting~---
objective or perceptual. In other words, what he considers more important: "what I
know" or "what I see." Imprecision in these matters leads to erroneous conclusions and
unconvincing constructions. Even so thoughtful a scholar as N.\,N.~Volkov, examining
in his recent monograph the question of "distortions" characteristic of works of art,
writes: "The depiction of parallel edges as converging, in linear perspective, is a
`distortion' (for in reality they are parallel)."\footnote{%
  N.\,N.~Volkov, \textit{Composition in Painting}. Moscow, 1977, p.~244.}
The error of such reasoning is connected with the fact that in perceptual space these
edges may well converge in linear perspective (if a person sees them so), and therefore
their representation in a picture need not be a "distortion." That they are in reality
parallel is important when depicting the other space~--- objective space, that is, in a
technical drawing. These two spaces and their depiction on a plane must not be
confused. Discussing a picture, one cannot appeal to the technical drawing; the geometric
properties of a picture and a technical drawing are different, but neither can be regarded
as a distortion of the other.

\srcpage{9}

The sculptor does not know all these difficulties. He conveys the true (objective)
geometry and, in so doing, the apparent geometry as well. Indeed, by conveying a
three-dimensional, solid object in three-dimensional sculpture, he enables one to form a
representation of its true geometry by examining the sculpture from different sides.
One need only pause and glance at the sculpture, and there arises in the human
consciousness a visual image of the object corresponding to that angle of view.

The difficulties that confront the painter~--- in contrast to the sculptor~--- are generated
by the geometric circumstance that he is forced to convey the appearance of
three-dimensional (solid) objects on a plane having only two dimensions. It is for this
reason that he stands before a choice: either objective geometry, or apparent geometry.
In contrast to sculpture, it is impossible to convey both simultaneously for all elements
of the image.

It is not unlikely that the reader will find this description of all the difficulties
connected with choosing between two possible geometries somewhat far-fetched. We
have long been accustomed to the fact that artists practically always convey visible
geometry, while engineers and draughtsmen deal with objective geometry. Yet this was
by no means always the case. Thinking abstractly, one must acknowledge that both
types of geometry have both strengths and weaknesses. If we turn to the history of art,
then at certain of its stages artists clearly preferred objective geometry to visible. The
reliefs and paintings of Ancient Egypt, for example, convey objective geometry. The
visual art of classical antiquity, by contrast, is clearly based on visible geometry.

One should not think that ancient Egyptian artists "could not" convey objects with
regard to perspective. They simply did not set themselves this task. If classical ancient art
were an unconditional and absolute step forward in the perfecting of pictorial means (in
Egypt they "could not," in Greece they "at last learned to"), then in medieval art there
would have been no revival of the aspiration to employ drafting methods. The
Renaissance, which once again completely "exterminated" drafting methods, was still
unable to defeat them "forever"; in our own time they are being born again in the newest
art. Consequently, the problem is by no means as simple and obvious as it might seem.
But here is not the place to consider questions connected with the stages of development
of the visual arts, with the causes that compelled artists repeatedly to change their view
of the character of the pictorial means best corresponding to artistic creation. It suffices
to establish that the problem facing the painter~--- objective or visible geometry?~--- cannot
be dismissed as far-fetched.

\srcpage{10}

The resolution of this problem in the artistic practice of peoples of different lands and
different epochs shows that in the vast majority of cases preference was and is given to
visible geometry. Many examples can be cited where painters convey only visible
geometry (for example, the paintings of the Renaissance); however, there also exist many
examples of the use in the visual arts of drafting methods (which permit the conveyance
of objective geometry), combined in one way or another with methods of pictorial
rendering (that is, with ways of conveying visible geometry). In the symbiosis of picture
and technical drawing, the leading role belongs, as a rule, to the picture; nevertheless
the techniques of technical drawing sometimes play an extremely important role. The
painter not infrequently attempts to convey both geometries simultaneously, in a single
work.

This combination is facilitated by the circumstance that the technical drawing, too, is
based on conveying visually perceivable properties of the depicted object, and in some
(rare) cases simply coincides with the picture. (To envisage this concretely, it suffices to
recall that a flat round object, which is normally perceived as an oval, can be seen as a
precise circle by choosing the appropriate viewpoint.)

A work of art may contain, alongside the geometry familiar to us, which is based on the
picture, the less familiar techniques of the technical drawing; and these elements by no
means indicate the artist's ineptitude or naivety. (This attribution is more likely to apply
to whoever comments on a work of art in this sense.) The painter in some cases attempts
to do what is easily accessible to the sculptor~--- to convey both visible and objective
geometry simultaneously. Of course, this proves possible only in part: some elements will
convey visible, others objective geometry; as already stated, it is simply impossible to
convey all elements of the image simultaneously in both geometries. Regardless of
whether the artist conveys only visible or only objective geometry, or combines them,
creating a kind of synthetic work, he will always communicate to the viewer only a
fraction of the geometric information that the sculptor is able to convey.

\srcpage{11}

There is yet another difficulty with which the painter inevitably contends. The fact is
that he lacks the means for the flawless conveyance of either objective or visible
geometry. Let us elucidate this by turning to the first of these tasks. Everyone who has
had occasion to deal with technical drawings knows that, as a rule, a drawing contains
several projections of the depicted object~--- for example, a front view, a side view, and
a top view. A technical drawing always presupposes a conditional depiction of the object,
and in the example cited the object is shown three times, from three conditional
directions of viewing. Such a threefold depiction of one and the same object leads to a
loss of immediacy, for in reality the viewer deals with one, not three, objects. Restoring
in the viewer's mind the true form of the object from three of its projections presupposes
analytical mental work, not the directly and immediately sensuous perception that fully
suffices when contemplating sculpture.

Moreover, the exact conveyance of visible geometry on the plane of a picture is also
fundamentally impossible. This problem will be examined in detail in subsequent
chapters; let us confine ourselves here to the observation that the depiction of the
visible appearance of nearby regions of space is possible only with "errors"~--- obvious
deviations from the visible forms of objects~--- and this is completely independent of the
type of perspectival construction chosen. The picture, like the technical drawing, is only
a conditional method of depicting space and the objects within it.

The painter not only cannot simultaneously convey both geometries (objective and
visible) for all that is depicted; even that which he decides to depict will contain
conventional elements, that is, distortions of directly and immediately sensuous
perception. In the search for an optimal solution to the artistic task confronting him, the
painter will seek the best ratio of elements of objective and visible geometry, and for
each of them he will be forced to introduce his own conventions~--- thereby opening
before the painter an immeasurably greater scope for varying geometric pictorial
means than those available to the sculptor (as already noted at the beginning of this
section, when the aspiration to convey the geometry of the depicted space and of the
objects within it with documentary precision is taken as a goal).

\srcpage{12}

The foregoing allows one to understand one peculiarity in the historical development of
sculpture in the round on the one hand, and of the depiction of spatial objects on a plane
(wall paintings, reliefs, canvases, and so on) on the other. If one compares sculpture in
the round from Ancient Egypt, Greece, and Rome, Gothic sculpture of Europe, sculpture
of the Renaissance, and the works of Rodin, then, although they convey entirely different
artistic images, all are sufficiently close to nature, and in terms of the pictorial means
employed they are in one way or another homogeneous. An ancient Egyptian sculptor,
seeing Michelangelo's \textit{David}, might have marvelled at the theme and the artistic
image, but he would have justly appreciated the plastic art of the master.

Quite a different picture emerges from a comparison of the works of painters of
different epochs. Here not only the artistic images are different, but the spatial pictorial
means will differ as well. In ancient Egyptian art there is a clear aspiration toward the
technical drawing; the Renaissance will consider as its ideal the exact following of the
visible geometry of the world; medieval art, in contrast to the art of the modern period,
employs reverse perspective; and so on. In doing this, painters of different epochs will
combine the shares of objective and visible geometry in their works in different ways and
will introduce differently chosen conventions. The types of spatial constructions on the
plane of the image will be different for different epochs and regions; and yet, as
paradoxical as it may seem, each of these types will be no more "correct" than another.
For the absolutely correct (fully corresponding to nature) can only be sculpture in the
round; in conveying space on a plane, the question of the "correctness" of the image
acquires meaning only if it is decided which geometric aspect of nature is to be conveyed
most exactly. This, in turn, correspondingly determines the choice of type of spatial
construction~--- which will of course change as soon as the orientation of the artist and
the viewer changes.

\srcpage{13}

One type of such construction will convey objective geometry more precisely, another
visible geometry, a third will combine them. These different types of spatial constructions
may be so far from one another that assessments of "better" and "worse" in comparing
different types of image sometimes become meaningless. Nevertheless, this need not
prevent an artist who employs one of the possible methods of spatial construction from
sincerely reproaching another artist~--- who prefers a different type~--- with a lack of
craftsmanship. Here one need only recall the opinions held of one another by
representatives of Chinese and European painting who "competed" at the court of the
last Chinese empress Cixi.

And yet sculptors of different epochs, from Ancient Egypt to Rodin, would have
understood each other better. Working in different historical epochs, they addressed
different artistic tasks in their sculptures; but the methods of spatial solution did not
stand in opposition to one another in sculpture in the round. On some notional scale of
absolute artistic values, the ancient Egyptian sculptural portrait of Nefertiti, the Gothic
sculpture of Uta from Naumburg Cathedral, and Michelangelo's Madonna from the
Medici Chapel would occupy approximately equal places~--- without any concession for
the supposed limitations of the art of earlier masters compared with later ones.

What is obvious for sculpture is, for some reason, not always considered valid in
assessing the work of painters. Too often, many people analysing, say, the works of the
masters of Ancient Egypt or the Middle Ages consciously or unconsciously compare them
with paintings of the Renaissance and with the canvases of modern painting, and
expressions of the following type are employed: "here the artist does not yet know how
to construct perspective correctly," "here the master seems to show the building in cross-
section." They usually refer to "dogmas" and "canons" constraining the artist's creativity,
seeing in this a justification for a certain "imperfection" of the works of art, and so on.

However, why should European painting of the sixteenth to nineteenth centuries be
regarded as the pinnacle of perfection, and why do these condescending intonations not
arise in the comparative consideration of sculpture from different epochs? Surely it was
not the case that talented and even brilliant people went into sculpture, while into
painting~--- before the Renaissance~--- only pitiful craftsmen went? Is it not more
reasonable to suppose that ancient Egyptian paintings, compared with the canvases of
the Renaissance, should have approximately the same comparative artistic value as the
sculpture of those same periods? But any modern viewer, entirely inexperienced in
questions of painting and sculpture, easily appreciates the works of ancient Egyptian
sculptors and regards ancient Egyptian paintings with a certain scepticism. This is
entirely natural: the pictorial means of the sculptors of past epochs in principle did not
differ from those in use today. In painting, what changed with time was not only its
artistic orientation and artistic image, but also the geometry of representation.

Here is not the place to seek an answer to the question of the causes of these changes;
but the very fact of the sharp change in the "geometry" of images is indisputable.
Analysis of the "geometries" employed by artists at different times shows that all are
purposeful, all are based on the real properties of the human psyche~--- in particular the
psychology of visual perception~--- all are in one way or another conditional, and the
"scientific" perspective of the Renaissance is by no means some absolute toward which
painters strove for centuries and with great difficulty. This "photographic" perspective
has been instilled in us since childhood through upbringing, and for this reason
deviations from it appear to many as inability or unwillingness to "draw correctly."
Yet this is not so. The domain in which the linear perspective familiar to us adequately
conveys visual perception is limited to the distant regions of space. As mathematical
analysis shows, for the near regions of space (when seeking to record exactly the visible
geometry of objects) one should employ axonometry and a gentle reverse perspective.
And if the painter wishes to convey not visible but objective geometry, he will be
compelled to employ techniques kindred to those that were in use in Ancient Egypt. The
techniques of the linear perspective of the Renaissance, of the medieval reverse
perspective, and of the special methods of ancient Egyptian depiction were the result of
the artist's position~--- conditioned by tradition, but consciously adopted.

\srcpage{14}

The history of the visual arts, and in particular the development of pictorial means,
must not be imagined as a kind of continuous and rectilinear ascent in the mastery of
new pictorial means, and comparative assessments of the perfection of works of art from
different epochs must not be built upon this image. But even confining ourselves to the
analysis of the techniques employed, we must know of the existence of different spatial
constructions. It is quite impossible to judge, for instance, the correctness and degree of
perfection of the geometric constructions characteristic of conveying space on a plane in
monuments of Ancient Egypt or the Middle Ages on the basis of the geometric dogmas
and canons of European painters of the sixteenth to nineteenth centuries. For in medieval
art and in the art of Ancient Egypt, different~--- but equally conditional and equally
geometrically grounded~--- spatial constructions were employed, just as in the painting of
the modern period. It would, after all, be absurd to evaluate ballet from the standpoint of
operatic art and to assert, for example, that "in ballet the actress already moves freely
about the stage, but has not yet learned to sing"~--- this is analogous to the reasoning
"has not yet learned to construct perspective."

The pictorial means were, of course, refined over time. But in analysing the process of
their development, one must distinguish between homogeneous methods of conveying
space on a plane~--- for example, those characteristic of classical antiquity and the
Renaissance~--- and different types of spatial constructions, for example those
characteristic of Ancient Egypt, which are foreign to the perspectival system of the
Renaissance.

Moreover, the question is legitimate: why did the methods of spatial constructions
change in precisely this (and not some other) sequence, why did art of the type of
ancient Egyptian painting precede the art of the Renaissance, and to what degree was
this course of development of pictorial means obligatory, conditioned by objective
causes? Why, having mastered the possibilities discovered in the Renaissance, do
painters today once again "discover" medieval and still earlier methods of conveying
space on the picture plane?

What interests us, however, is the objective properties of various methods of depicting
space on a plane, their conditionality by the regularities of visual perception and other
rational circumstances, and the analysis of their comparative possibilities.

The frequent comparisons throughout the book between the art of Ancient Egypt, the
Middle Ages, and the Renaissance are prompted by a desire to lend the geometric
arguments a certain vividness, and to remind the reader of the most familiar monuments
of art. If one examines how geometric pictorial means were applied in the art of various
peoples in depicting space on a plane, it turns out that the geometric distinctiveness of
these means is connected, first, with which geometry~--- visible or objective~--- the artist
and viewer considered the more important, and second, with the conventions adopted in
depicting it (since neither the one nor the other can in principle be conveyed absolutely
flawlessly, without resort to more or less artificial conventions).

All that has been said applies not only to past epochs but is relevant today and will
remain relevant in the future. For the problem of the "documentary-precise" conveyance
of spatial images on a plane (we have in mind the geometric side of the problem) is an
"eternal" one, since it has no unambiguous mathematical solution. A departure from the
perspectival techniques of modern art by no means signifies a departure from realism;
it may be a transition to other, mathematically equally well-grounded methods of
conveying space on a plane. Of course, documentary-precise conveyance of space is not,
as already noted, the artist's goal; however, an analysis, based on mathematics and the
psychology of perception, of the methods of its realization opens up the possibility of
assessing the potential possibilities contained in this or that method of spatial
construction, and of understanding where and to what degree the artist departs from
mathematically grounded schemes.

What for the sculptor presents no problem whatsoever turns out for the painter to be a
complex task with no unambiguous mathematical solution. This, on the one hand,
complicates the painter's work; on the other, it increases the diversity of pictorial means
and correspondingly broadens his possibilities.

Below, a number of typical systems of spatial constructions characteristic of different
epochs and regions will be examined on the basis of the laws of mathematics and the
peculiarities of human perception.
